\documentstyle[%


righttag%


,11pt%


,amssymb%


,russian%               remove in Eng.


%,izvru%                 Set izven% in Eng.


,izvru-ks%             for brief communications


]{amsart}





%       Math definitions





\newcommand{\Int}{\operatorname{Int}}


\newcommand{\tts}[1]{}





%\hoffset=-15mm%                  For Eng.


\setcounter{page}{62}


\god{2001}


\nomer{4~(467)} %                        Remove in Eng.


%\nomer{Vol.\,45, No.\,4}%               Set in Eng.


%\str{\pageref{begin}--\pageref{end}}%  Set in Eng.


\udk{517.982}





\author{\it Л.В.\,Веселова, О.Е.\,Тихонов}


% NB:   ^^ remove in Eng.


\title{O единственности решения обратных задач


интерполяции для классов операторов}


\thanks{Работа выполнена при поддержке Российского фонда


фундаментальных исследований, грант \No\,98-01-00103.}





\date{}


%\pagestyle{myheadings}%         Set in Eng.


\pagestyle{plain} %              Remove in Eng.


\begin{document}


%\thispagestyle{plain}%          Set in Eng.


\maketitle


\label{begin}





В работах [1], [2] была доказана единственность решения обратных 


задач интерполяции линейных операторов в классической постановке (см.


также [3]). Цель данной работы --- обобщение этих 


результатов в рамках достаточно общего подхода к интерполяции для классов 


операторов (в том числе --- нелинейных), а также их аналог для 


интерполяции положительных линейных операторов в банаховых решетках. 


\medskip





{\bf 1. Обозначения} (по поводу определений см., напр., [4], [5]).


Пусть $X$ --- банахово пространство,


$\overline X$ --- банахова пара.


Через $B(X)$ и $L(X)$ (соответственно $B(\overline X)$ и $L(\overline X)$)


обозначаем банахово пространство всех ограниченных нелинейных операторов и


подпространство ограниченных линейных операторов в $X$ (в


$\overline X$).  


Подмножества $L(X)$ и $L(\overline X)$, образованные операторами


единичного ранга, обозначаем $R1(X)$ и $R1(\overline X)$


соответственно.





Для двух банаховых пространств $X$ и $Y$ запись $X \hookrightarrow Y$ 


означает, что $X$ линейно и топологически вложено в $Y$, 


запись $X \simeq Y$ --- что $X$ и $Y$ совпадают как 


линейные пространства и нормы в них эквивалентны; если нормы 


пропорциональны, то пишем $X \cong Y$.





Для двух банаховых пар $\overline X = (X_0,X_1)$ и 


$\overline Y = (Y_0,Y_1)$ запись $\overline X \simeq \overline Y$ 


означает, что либо $X_i \simeq Y_i$,  


либо $X_i \simeq Y_{|i-1|}$ ($i=0,1$), 


запись $\overline X \cong \overline Y$ --- 


что либо $X_i \cong Y_i$, либо $X_i \cong Y_{|i-1|}$ ($i=0,1$). 


\medskip





{\bf 2. Схема интерполяции для классов операторов.}


С каждой банаховой парой $\overline X =(X_0,X_1)$ считаем ассоциированым


некоторый класс $\kappa (\overline X)$ и некоторое банахово пространство


$\sigma (\overline X)$ операторов в $\overline X$, причем


$R1(\overline X) \subset \kappa (\overline X) \subset


\sigma (\overline X) \hookrightarrow B(\overline X)$ и


$\| S \| _{\sigma (\overline X)}= \| S \| _{L(\overline X)}$ для


любого $S \in R1(\overline X)$.


С каждым банаховым пространством $X$ считаем ассоциированым класс


операторов $\lambda (X)$ и банахово пространство $\tau (X)$, причем 


$R1(X) \subset \lambda (X) \subset \tau (X) \hookrightarrow B(X)$,


$\| S \| _{\tau (X)} = \| S \| _{L(X)}$ для любого $S \in R1(X)$ и,


кроме того, для любого $T \in \kappa (\overline X)$


ограничение $T | _{X_i}$ принадлежит $\lambda (X_i)$ и


$\| T | _{X_i} \| _{ \tau (X_i)} \le


\| T \| _{\sigma (\overline X)}$ $(i=0,1)$.





\begin{defn}


Пусть $\frak I=(\kappa , \sigma , \lambda , \tau )$ ---


конкретная реализация вышеприведенной схемы.


Промежуточное для банаховой пары $\overline X$ пространство $Z$


называем $\frak I${\it-интерполяционным}, если существует константа 


$c>0$ такая, что $T | _Z \in \lambda (Z)$ и


$\| T | _Z \| _{ \tau (Z)} \le c \| T \| _{\sigma (\overline X)}$


для любого $T \in \kappa (\overline X)$. Если можно взять $c=1$, то такое


пространство $Z$ называем {\it нормально $\frak I$-интерполяционным}.


\end{defn}





Совокупность всех $\frak I$-интерполяционных пространств для банаховой 


пары $\overline{X}$ обозначаем далее $\frak I$-$\Int\overline X$,


совокупность всех нормально $\frak I$-интерполяционных пространств 


--- $\frak I$-$\Int_1 \overline X$.





Заметим, что при


$\kappa (\overline X)=\sigma (\overline X)=L(\overline X)$,


$\lambda (X)=\tau (X)=L(X)$ приходим к обычному понятию 


интерполяционности относительно ограниченных линейных операторов. 


Соответствующую совокупность $\frak I$-$\Int\overline X$  


обозначаем далее $\Int\overline X$.





Взяв $\kappa (\overline X)=R1(\overline X)$,


$\sigma (\overline X)=L(\overline X)$, $\lambda (X)=R1(X)$, $\tau (X)=L(X)$,


получаем известное определение интерполяционности для операторов единичного 


ранга. 


Соответствующие совокупности интерполяционных пространств 


обозначаем далее 


$R1$-$\Int\overline X$ и $R1$-$\Int_1 \overline X$.





Ряд других содержательных примеров получается при изучении интерполяции 


идеалов операторов в банаховых парах.


Проиллюстрируем естественность приведенной схемы на примере интерполяции 


ядерных операторов. Пусть известно, что изучаемые операторы являются 


ядерными как в пространстве $X_0$, так и в $X_1$. Интересны


следующие вопросы для промежуточного пространства $Z$.





а) Переводят ли рассматриваемые операторы $Z$ в $Z$ и справедливы ли 


интерполяционные оценки для обычных операторных норм?





б) Являются ли вдобавок эти операторы ядерными как операторы из $Z$ в 


$Z$?





в) Можно ли получить интерполяционные оценки для ядерных норм?





Каждой из постановок соответствует своя реализация вышеприведенной схемы. 


Например, для задачи~б) естественно взять 


${\cal N}(X_0) \cap {\cal N}(X_1)$ в качестве $\kappa (\overline X)$, 


$L(\overline X)$ --- в качестве $\sigma (\overline X)$, 


${\cal N}(X)$ --- в качестве $\lambda (X)$, $L(X)$ --- в качестве $\tau (X)$ 


(здесь ${\cal N}(X)$ --- совокупность ядерных операторов в банаховом 


пространстве $X$). Как следует из примера Пича (см. [6], п.\,1.19.8), 


$\frak I$-$\Int\overline X$ при этом подходе может не содержать 


$\Int\overline X$.





Для некоторых пар $\frak I$-интерполяционные пространства различных схем 


$\frak I$ могут совпадать, однако, как правило, различные конкретные 


реализации общей схемы дают, вообще говоря, 


различные совокупности $\frak I$-интерполяционных пространств.  


\medskip





{\bf 3. Единственность решения обратных задач интерполяции.}





\begin{thm}


Пусть $\frak I$ --- некоторая реализация схемы интерполяции 


классов операторов п.\,$2$ и пусть для двух банаховых пар $\overline X$ 


и $\overline Y$ выполнено соотношение


$\frak I$-$\Int\overline X =


\frak I$-$\Int\overline Y$.


Тогда $\overline X \simeq \overline Y$.


\end{thm}





Для доказательства этой теоремы достаточно заметить, что для банаховых 


пар $\overline X = (X_0,X_1)$ и $\overline Y = (Y_0,Y_1)$


соотношение 


$\frak I\text{-}\Int \overline X =


\frak I\text{-}\Int\overline Y$ 


влечет 


$X_i \in R1\text{-}\Int \overline Y$ и


$Y_i \in R1\text{-}\Int \overline X$ ($i=0,1$), 


и применить следующую лемму. 





\begin{lem}


Если для двух банаховых пар $\overline X =(X_0,X_1)$ и


$\overline Y = (Y_0,Y_1)$ выполняются условия


$X_i \in R1\text{-}\Int \overline Y$, 


$Y_i \in R1\text{-}\Int \overline X$ $(i=0,1)$, 


то $\overline X \simeq \overline Y$.


\end{lem}





В справедливости этой леммы можно убедиться, проанализировав 


доказательства работы [1]. 


Заметим еще, что из результатов работы [3] (предложение 1.4 и теорема 2.8) 


лемма 1 следует непосредственно.





Аналогично тому, как теорема 1 сводится к лемме 1, следующая теорема 2 


сводится к соответствующей лемме 2.





\begin{thm}


Пусть для двух банаховых пар $\overline X$ и $\overline Y$ выполнено 


соотношение


$\frak I\text{-}\Int_1 \overline X =


\frak I\text{-}\Int_1 \overline Y$.


Тогда $\overline X \cong \overline Y$.


\end{thm}





\begin{lem}


Если для двух банаховых пар $\overline X =(X_0,X_1)$ и


$\overline Y = (Y_0,Y_1)$ выполняются условия


$X_i \in R1\text{-}\Int_1 \overline Y$, 


$Y_i \in R1\text{-}\Int_1 \overline X$ $(i=0,1)$, 


то $\overline X \cong \overline Y$.


\end{lem}





В справедливости леммы 2 можно убедиться, проанализировав доказательства 


работы [2], принимая во внимание лемму 1. 


Как и лемма 1, из результатов работы [3] (предложение 1.4 и 


теорема 3.7) лемма 2 следует непосредственно.


\medskip





{\bf 4. Единственность решения обратных задач интерполяции положительных 


oператоров в банаховых решетках.}


Напомним





\begin{defn}[{\series{m}\shape{n}\selectfont [7]}]


Говорят, что две банаховы решетки $X_0$ и $X_1$ образуют {\it 


интерполяционную пару банаховых решеток,} 


если они обе алгебраически и топологически вложены в некоторое 


отделимое топологическое векторное пространство, которое к тому же 


является векторной решеткой, причем и $X_0$, и $X_1$ --- идеалы в этой 


решетке.     


\end{defn}





Отметим, что в рассматриваемой ситуации банаховы пространства 


пересечения $X_0 \cap X_1$ и суммы $X_0 + X_1$ в свою очередь оказываются 


банаховыми решетками. 





Для двух банаховых решеток $E$ и $F$ запись $E \hookrightarrow F$ далее 


дополнительно подразумевает, что $E$ --- идеал в $F$.





\begin{defn}


Пусть $\overline X = (X_0, X_1)$ --- интерполяционная пара банаховых 


решеток. Банахова решетка $Z$, удовлетворяющая условию


$X_0 \cap X_1 \hookrightarrow Z \hookrightarrow X_0 + X_1$, 


называется {\it положительно интерполяционной}, если существует константа 


$c>0$ такая, что для любого линейного положительного оператора $T$, 


действующего в $\overline X$, имеем


$T(Z) \subset Z$ и $\| T \| _{Z \to Z} \le c \| T \| _{L(\overline X)}.$


\end{defn}





Совокупность положительно интерполяционных решеток обозначим 


$P\text{-}\Int \overline X$, нормально положительно интерполяционных --- 


$P\text{-}\Int_1 \overline X$. 





Положительная интерполяционность изучалась, в основном, в связи с 


конструкцией Каль\-де\-ро\-на--Лозановского.  


Отметим, что, вообще говоря, положительно интерполяционные решетки не 


обязательно интерполяционны для всех линейных ограниченных операторов 


[8].





\begin{lem}


Для любой интерполяционной пары $\overline X$ банаховых решеток 


выполняются соотношения 


$P\text{-}\Int \overline X \subset R1\text{-}\Int \overline X$ и 


$P\text{-}\Int_1 \overline X \subset R1\text{-}\Int_1 \overline X$.


\end{lem} 





В следующей теореме использование символов ``$\simeq $'' и ``$\cong $''  


подразумевает, вдобавок к эквивалентности или 


пропорциональности норм, и совпадение отношения порядка на 


соответствующих банаховых решетках.





\begin{thm}


Пусть $\overline X$, $\overline Y$ --- две интерполяционные пары


банаховых решеток.





{\em a)} Если 


$P\text{-}\Int \overline X = P\text{-}\Int \overline Y$, то 


$\overline X \simeq \overline Y$.





{\em b)} Если 


$P\text{-}\Int_1 \overline X = P\text{-}\Int_1 \overline Y$, то 


$\overline X \cong \overline Y$.


\end{thm}





Доказательство этой теоремы сводится к применению лемм 1, 2 и 3.





\begin{thebibliography}{9}





\bibitem{1}


Tikhonov O.E., Veselova L.V. {\it A Banach couple is determined 


by the collection of its interpolation spaces} // Proc.\ Amer.\ 


Math.\ Soc. -- 1998. -- V.\,126. -- \No\,4. -- P.\,1049--1054. 


\bibitem{2}


Tikhonov O.E., Veselova L.V. {\it The uniqueness of the solution 


to the inverse problem of exact interpolation} // Israel Math.\ 


Conf.\ Proc. -- 1999. -- V.\,13. -- P.\,209--215.


\bibitem{3}


Веселова Л.В., Тихонов О.Е. {\it О единственности решения


обратных задач интерполяции} // Препринт НИИММ \No\,95-2. -- Казань: Казанск.


фонд ``Математика'', 1995. -- 17 с.


% (Англ. перевод: XXX e-Print 


% Archive, math.FA/9902108).


\bibitem{4}


Брудный Ю.А., Крейн С.Г., Семенов Е.М. {\it Интерполяция 


линейных операторов} // Итоги науки и техн. Матем.\ анализ.


-- М.: ВИНИТИ, 1986. -- T.\,24. -- С.\,3--163. 


\bibitem{5}


Ovchinnikov V.I. {\it The method of orbits in interpolation theory} // 


Math.\ Rept. -- 1984. -- V.\,1. -- P.\,349--515.  


\bibitem{6}


Трибель Х. {\it Теория интерполяции, функциональные пространства,


дифференциальные операторы.} -- М.: Мир, 1980. -- 664 c.


\bibitem{7}


Lindenstrauss J., Tzafriri L. {\it Classical Banach spaces}. II. 


{\it Function spaces.} -- Berlin--Heidelberg--New York: Springer-Verlag, 1979.


-- 243 p.


\bibitem{8}


Лозановский Г.Я. {\it Замечание об одной интерполяционной теореме 


Кальдерона} // Функц.\ анализ и его прилож. -- 1972. -- Т.\,6. -- Вып.\,6. 


-- С.\,89--90.





\end{thebibliography}





\vskip 0.5cm





\post{\it Казанский государственный}{\it Поступила}


\post{\it технологический университет}{\it $20.03.2000$}


\smallskip


\post{\it НИИ математики и механики при}{}


\post{\it Казанском государственном университете}{}





\label{end}


\end{document}


